% Informe del sprint 1 del framework de desarrollo seguro.
% Compilar con XeLaTeX (dos pasadas para el índice):
%   latexmk -xelatex informe_sprint1.tex
\documentclass[11pt,a4paper]{report}

% ---------- Idioma y tipografía ----------
\usepackage{fontspec}
\usepackage[spanish,es-noshorthands,es-tabla]{babel}
\setmonofont{Latin Modern Mono}[Scale=0.95]

% ---------- Página ----------
\usepackage[a4paper,top=2.6cm,bottom=2.6cm,left=2.5cm,right=2.5cm,headheight=14pt]{geometry}
\usepackage{xcolor}
\usepackage{titlesec}
\usepackage{fancyhdr}
\usepackage{booktabs}
\usepackage{tabularx}
\usepackage{longtable}
\usepackage{array}
\usepackage{colortbl}
\usepackage{enumitem}
\usepackage{caption}
\usepackage{float}
\usepackage[most]{tcolorbox}
\usepackage{tikz}
\usetikzlibrary{arrows.meta,positioning,fit,backgrounds}
\usepackage{xurl}
\usepackage[hidelinks]{hyperref}

% ---------- Colores (los mismos del plan técnico v2) ----------
\definecolor{azulfw}{HTML}{2A64A8}
\definecolor{azulclaro}{HTML}{E4EEF9}
\definecolor{verdefw}{HTML}{3E6B12}
\definecolor{verdeclaro}{HTML}{EAF1E1}
\definecolor{ocrefw}{HTML}{8A4E0B}
\definecolor{ocreclaro}{HTML}{FCF0DF}
\definecolor{rojofw}{HTML}{A33A3A}
\definecolor{rojoclaro}{HTML}{F8E3E3}
\definecolor{grisfw}{HTML}{6B6B6B}

% ---------- Títulos ----------
\titleformat{\chapter}[hang]{\huge\bfseries\color{azulfw}}{\thechapter.}{0.5em}{}
\titlespacing*{\chapter}{0pt}{0pt}{1.2em}
\titleformat{\section}{\Large\bfseries\color{azulfw}}{\thesection}{0.6em}{}
\titleformat{\subsection}{\large\bfseries}{\thesubsection}{0.6em}{}
\addto\captionsspanish{\renewcommand{\contentsname}{\color{azulfw}Índice}}

% ---------- Encabezado y pie ----------
\pagestyle{fancy}
\fancyhf{}
\fancyhead[L]{\small\color{grisfw}Framework de desarrollo seguro guiado por conocimiento}
\fancyhead[R]{\small\color{grisfw}Informe del sprint 1}
\fancyfoot[C]{\small\thepage}
\renewcommand{\headrulewidth}{0.4pt}
\fancypagestyle{plain}{\fancyhf{}\fancyfoot[C]{\small\thepage}\renewcommand{\headrulewidth}{0pt}}

% ---------- Tablas ----------
\renewcommand{\arraystretch}{1.25}
\newcolumntype{L}{>{\raggedright\arraybackslash}X}
\newcolumntype{P}[1]{>{\raggedright\arraybackslash}p{#1}}
\captionsetup{font=small,labelsep=colon}
\setlist{itemsep=0.25em,topsep=0.4em}

% ---------- Recuadros ----------
\newtcolorbox{cajaverde}[1]{enhanced,colback=verdeclaro,colframe=verdefw,
  coltitle=white,fonttitle=\bfseries,title={#1},arc=1.5mm,boxrule=0.6pt,
  left=2mm,right=2mm,top=1mm,bottom=1mm}
\newtcolorbox{cajaazul}[1]{enhanced,colback=azulclaro,colframe=azulfw,
  coltitle=white,fonttitle=\bfseries,title={#1},arc=1.5mm,boxrule=0.6pt,
  left=2mm,right=2mm,top=1mm,bottom=1mm}
\newtcolorbox{cajaocre}[1]{enhanced,colback=ocreclaro,colframe=ocrefw,
  coltitle=white,fonttitle=\bfseries,title={#1},arc=1.5mm,boxrule=0.6pt,
  left=2mm,right=2mm,top=1mm,bottom=1mm}

\newcommand{\pr}[1]{\href{https://github.com/PracticasESPE/Framework_Desarrollo_Seguro/pull/#1}{\##1}}
\newcommand{\filasprint}[1]{\rowcolor{azulclaro}\multicolumn{5}{l}{\bfseries #1}\\}

\hypersetup{pdftitle={Informe del sprint 1},pdfauthor={Marcelo Acuña Gamboa, Abner Arboleda}}

\begin{document}

% ======================================================================
% PORTADA
% ======================================================================
\begin{titlepage}
\centering
\vspace*{4.2cm}
{\color{azulfw}\rule{\textwidth}{0.8pt}}\par\vspace{0.9cm}
{\fontsize{24}{30}\selectfont\bfseries Framework de desarrollo seguro\\
basado en métodos computacionales\\ guiados por conocimiento\par}
\vspace{0.9cm}{\color{azulfw}\rule{\textwidth}{0.8pt}}\par
\vspace{1.1cm}
{\Large Informe del sprint 1\par}
\vspace{0.3cm}
{\large Walking skeleton y fase 4 inicial (Security Gate 1)\par}
\vspace{0.3cm}
{\normalsize 5 al 16 de octubre de 2026\par}
\vspace{2.2cm}
{\normalsize
\begin{tabular}{rl}
Marcelo Acuña Gamboa & Participante 1: motor e IA\\
Abner Arboleda & Participante 2: plataforma y seguridad\\
\end{tabular}\par}
\vfill
{\normalsize Universidad de las Fuerzas Armadas – ESPE \quad$\bullet$\quad Ingeniería de Software\par}
\vspace{0.4cm}
{\normalsize Versión 1.0 \quad$\bullet$\quad 16 de octubre de 2026\par}
\end{titlepage}

\tableofcontents

% ======================================================================
\chapter{Resumen del sprint}
% ======================================================================

\section{Objetivo e hito}

El sprint 1 corresponde al principio P3 del plan técnico, \emph{primero de punta a punta,
después inteligente}. Su objetivo era construir las herramientas del Security Gate 1 y hacer
que un microservicio NestJS escrito a mano lo superara en GitHub Actions, con sus hallazgos
registrados en Neo4j. También había que elegir el proveedor de LLM.

\begin{cajaazul}{Hito de control del 16 de octubre de 2026}
\emph{Fase 4 inicial: un microservicio escrito a mano pasa el Security Gate 1 en CI. Si no se
cumple, no se avanza a la generación con LLM.}
\end{cajaazul}

\begin{cajaverde}{Resultado: hito cumplido}
El servicio de referencia \texttt{usuarios} superó el Security Gate 1 en la ejecución
\href{https://github.com/PracticasESPE/Framework_Desarrollo_Seguro/actions/runs/38955751601}{38955751601}
de GitHub Actions sobre la rama \texttt{dev} (16/10/2026). Las ocho comprobaciones terminaron en
\emph{OK} al primer intento y los cinco archivos SARIF se publicaron en la pestaña Security. La
fase 4 del motor registró la ejecución en Neo4j. Se completaron todas las tareas planificadas
del sprint y se puede pasar al sprint 2.
\end{cajaverde}

\section{Indicadores}

\begin{table}[H]
\centering
\caption{Indicadores del sprint 1 medidos en la última ejecución sobre \texttt{dev}.}
\begin{tabularx}{\textwidth}{@{}P{6.2cm}L@{}}
\toprule
\textbf{Indicador} & \textbf{Valor}\\
\midrule
Días y horas del sprint & 9 días, 72 horas de equipo (36 por integrante)\\
Tareas planificadas completadas & 18 de 18 (9 por integrante)\\
Pull requests fusionados en \texttt{dev} & 17 (\pr{137} a \pr{155}), todos con revisión del otro integrante\\
Duración del pipeline completo & 3 min 25 s (Security Gate 1: 46 s)\\
Resultado del Security Gate 1 & 8 de 8 comprobaciones en \emph{OK}, 0 hallazgos\\
Reglas Semgrep aplicadas & 82: 74 públicas (\texttt{p/typescript}, \texttt{p/nodejs}) y 8 propias\\
Casos de prueba de las reglas propias & 23 detecciones esperadas, 23 obtenidas; 0 falsos positivos en los casos \emph{ok}\\
Pruebas del motor (pytest) & 106 pruebas en verde, 95\,\% de cobertura de \texttt{orchestrator/}\\
Pruebas del servicio (Jest) & 16 pruebas en verde, 80,8\,\% de cobertura de sentencias\\
Esquema de Neo4j & 14 restricciones y 7 índices, aplicados en cada ejecución de CI\\
Proveedor de LLM & Groq (\texttt{openai/gpt-oss-20b} y \texttt{gpt-oss-120b})\\
Alertas abiertas en GitHub Security & 0\\
\bottomrule
\end{tabularx}
\end{table}

% ======================================================================
\chapter{Trabajo realizado}
% ======================================================================

\section{Cumplimiento del cronograma}

La tabla compara el cronograma diario del plan con lo entregado. La columna PR indica el pull
request con el que cada tarea entró en \texttt{dev}. Las filas de un solo bloque son trabajo conjunto.

{\small
\begin{longtable}{@{}P{1.55cm}P{4.6cm}P{4.6cm}P{1.6cm}P{1.6cm}@{}}
\caption{Tareas del sprint 1 por integrante y pull request asociado.}\\
\toprule
\textbf{Fecha} & \textbf{Marcelo (motor e IA)} & \textbf{Abner (plataforma y seguridad)} & \textbf{PR Marcelo} & \textbf{PR Abner}\\
\midrule
\endfirsthead
\toprule
\textbf{Fecha} & \textbf{Marcelo (motor e IA)} & \textbf{Abner (plataforma y seguridad)} & \textbf{PR Marcelo} & \textbf{PR Abner}\\
\midrule
\endhead
\filasprint{Sprint 1: Walking skeleton y fase 4 inicial (Security Gate 1)}
05/10 lun & Esqueleto del motor: uv, Typer, FastAPI y configuración. & \texttt{docker-compose.dev.yml} con Neo4j, RabbitMQ y Keycloak. & \pr{139} & \pr{137}\\
06/10 mar & Interfaz \texttt{ProveedorLLM} y adaptadores (Claude, Gemini, Ollama). & Servicio de referencia \texttt{usuarios}: módulo, DTO con class-validator y Swagger. & \pr{140} & \pr{138}\\
07/10 mié & Prueba comparativa con 5 requisitos; elección de Groq. & Prisma + PostgreSQL y pruebas Jest base. & \pr{141} & \pr{142}\\
08/10 jue & Modelo común de hallazgos y lector SARIF. & Dockerfile multietapa con Hadolint; workflow con un paso por fase. & \pr{144} & \pr{143}\\
12/10 lun & Esquema de Neo4j: restricciones e índices en \texttt{graph/}. & Gate 1 en CI: tsc, ESLint, Semgrep, Gitleaks, Trivy y Hadolint. & \pr{148} & \pr{146}\\
13/10 mar & Registro de ejecuciones y hallazgos como nodos en Neo4j. & SARIF en la pestaña Security; \texttt{npm ci -{}-ignore-scripts} y \texttt{.npmrc}. & \pr{151} & \pr{147}\\
14/10 mié & Comando \texttt{framework run -{}-fase 4} integrado al workflow. & Reglas Semgrep propias v1 para NestJS y Prisma. & \pr{152}, \pr{153} & \pr{149}\\
15/10 jue & La fase 4 lee los SARIF del Gate 1 y los registra en Neo4j. & Gate 1 estable en el runner: reintentos, resumen y caché de Trivy. & \pr{155} & \pr{150}\\
16/10 vie & Pruebas del motor con pytest y guía de demostración. & Correcciones del pipeline, este informe y planificación del sprint 2. & \pr{155} & este PR\\
\bottomrule
\end{longtable}}

Algunas tareas se fusionaron uno o dos días después de su fecha planificada; por ejemplo,
\pr{143} se fusionó el 11/10 y \pr{151} el 16/10. Estos retrasos se absorbieron dentro del
sprint y no afectaron al hito.

\section{Plataforma y seguridad}

\subsection{Servicio de referencia \texttt{usuarios}}
Es un microservicio NestJS escrito a mano que sirve de banco de pruebas para todas las fases
antes de que exista la generación con LLM. Expone tres operaciones REST (crear, listar y
consultar por identificador) documentadas con Swagger. Persiste los datos en PostgreSQL mediante
Prisma y valida la entrada con un \texttt{ValidationPipe} global en modo \texttt{whitelist}. Las
consultas de lista tienen un límite de 100 registros. Las 16 pruebas Jest cubren el DTO, el
controlador y el servicio.

\subsection{Contenerización}
El Dockerfile tiene tres etapas: compilación, dependencias de producción e imagen final. La imagen
final contiene solo el JavaScript compilado y las dependencias de producción, se ejecuta con un
usuario sin privilegios (\texttt{USER 1000:1000}) e incluye \texttt{HEALTHCHECK}. Hadolint no
reporta ningún aviso, ni siquiera de nivel \emph{info}.

\subsection{Security Gate 1}
El Gate 1 se ejecuta en la fase 4 del workflow mediante \texttt{scripts/gate1.sh}. Aplica ocho
comprobaciones al servicio de referencia y genera un SARIF por herramienta.

\begin{table}[H]
\centering
\caption{Comprobaciones del Security Gate 1 y resultado en la ejecución del hito.}
\begin{tabularx}{\textwidth}{@{}P{3.3cm}LP{2.4cm}P{2.0cm}@{}}
\toprule
\textbf{Herramienta} & \textbf{Qué busca} & \textbf{SARIF} & \textbf{Resultado}\\
\midrule
\texttt{tsc -{}-noEmit} & Compilación y tipos & -- & OK\\
ESLint + security & Patrones inseguros de JavaScript y estilo & -- & OK\\
Semgrep & SAST con 74 reglas públicas y 8 propias & \texttt{semgrep} & OK\\
\texttt{semgrep -{}-test} & Que las reglas propias detecten sus casos de prueba & -- & OK\\
Gitleaks & Secretos en todo el historial de Git & \texttt{gitleaks} & OK\\
Trivy fs & Vulnerabilidades HIGH y CRITICAL en \texttt{package-lock.json} y \texttt{uv.lock} & \texttt{trivy-fs} & OK\\
Trivy config & Configuraciones inseguras en el Dockerfile & \texttt{trivy-config} & OK\\
Hadolint & Buenas prácticas del Dockerfile & \texttt{hadolint} & OK\\
\bottomrule
\end{tabularx}
\end{table}

El script distingue entre \emph{hallazgos}, cuando la herramienta funciona y encuentra problemas,
y \emph{errores de la herramienta}, como fallos de red o de descarga. Los hallazgos ponen el Gate
en rojo sin reintentar. Los errores se reintentan hasta dos veces, para que un fallo transitorio
del runner no se confunda con un problema de seguridad. Al final, el script escribe una tabla de
resultados en el resumen de la ejecución.

\subsection{Reglas Semgrep propias (v1)}
Las reglas de \texttt{semgrep\_rules/} forman parte de la base de conocimiento (principio P5). Cada
una declara el CWE y la categoría OWASP que detecta, para que el cargador del sprint 2 pueda
enlazarlas en Neo4j. El archivo \texttt{nestjs-seguridad.ts} contiene los casos de prueba, y el
pipeline comprueba en cada ejecución que todas las reglas siguen detectándolos.

\begin{table}[H]
\centering
\caption{Reglas Semgrep propias, versión 1.}
\begin{tabularx}{\textwidth}{@{}>{\raggedright\arraybackslash\footnotesize}p{5.65cm}LP{2.0cm}P{1.1cm}@{}}
\toprule
\textbf{Regla} & \textbf{Qué detecta} & \textbf{Severidad} & \textbf{CWE}\\
\midrule
\texttt{nestjs-validationpipe-sin-whitelist} & Asignación masiva por propiedades no declaradas & WARNING & 915\\
\texttt{prisma-sql-crudo-inseguro} & \texttt{\$queryRawUnsafe} y \texttt{\$executeRawUnsafe} & ERROR & 89\\
\texttt{nestjs-cors-abierto} & CORS abierto a cualquier origen & WARNING & 942\\
\texttt{jwt-secreto-en-codigo} & Secreto JWT escrito en el código & ERROR & 798\\
\texttt{jwt-ignora-expiracion} & Tokens vencidos aceptados & ERROR & 613\\
\texttt{tls-verificacion-desactivada} & Certificados TLS sin verificar & ERROR & 295\\
\texttt{hash-criptografico-debil} & MD5 o SHA-1 & WARNING & 327\\
\texttt{prisma-findmany-sin-limite} & Consultas de lista sin \texttt{take} & WARNING & 770\\
\bottomrule
\end{tabularx}
\end{table}

\section{Motor e IA}

\begin{itemize}
\item \textbf{Esqueleto del motor} (\pr{139}): proyecto Python con uv, CLI \texttt{framework} con
Typer, API con FastAPI y configuración por variables de entorno.
\item \textbf{Adaptador de LLM} (\pr{140}, \pr{141}): interfaz común con salida estructurada en
Pydantic y adaptadores para Claude, Gemini, Ollama y Groq.
\item \textbf{Modelo de hallazgos y lector SARIF} (\pr{144}): convierte la salida de cualquier
herramienta al modelo común, como pide el principio P2. Se probó con ejemplos y con salidas
reales de Hadolint y Gitleaks.
\item \textbf{Neo4j} (\pr{148}, \pr{151}): 14 restricciones y 7 índices, y registro de nodos
\texttt{Ejecucion} y \texttt{Hallazgo} enlazados con su \texttt{CWE}.
\item \textbf{Fase 4 en CI} (\pr{152}, \pr{153}, \pr{155}): el workflow levanta un Neo4j propio
de la ejecución, ejecuta las pruebas del motor, aplica el esquema y registra los hallazgos del
Gate 1. Los SARIF de ejemplo sirven para mostrar cómo se ve el grafo cuando sí hay hallazgos.
\end{itemize}

\begin{table}[H]
\centering
\caption{Comparación de proveedores de LLM con cinco requisitos (7/10/2026).}
\begin{tabularx}{\textwidth}{@{}LP{2.6cm}P{2.6cm}@{}}
\toprule
\textbf{Criterio} & \textbf{Gemini} & \textbf{Groq}\\
\midrule
JSON válido al primer intento & 5/5 & 5/5\\
Clasificación RF/RNF correcta & 5/5 & 5/5\\
Cobertura de etiquetas de seguridad & 0,80 & 0,77\\
TypeScript generado / que compila & 1/5 / 1 & 5/5 / 5\\
Latencia media de generación & 110,3 s & 2,0 s\\
Reintentos por saturación o cuota & 10 & 0\\
\bottomrule
\end{tabularx}
\end{table}

\begin{cajaazul}{Decisión: Groq como proveedor de LLM}
Las dos opciones extraen requisitos con la misma calidad. Groq, en cambio, generó y compiló todo
el TypeScript sin reintentos y fue unas 55 veces más rápido en generación. Gemini agotó la cuota
gratuita durante la prueba. El adaptador permite cambiar de proveedor sin tocar el resto del motor.
\end{cajaazul}

% ======================================================================
\chapter{Estado del pipeline al cierre}
% ======================================================================

\section{Walking skeleton hasta el Gate 1}

\begin{figure}[H]
\centering
\begin{tikzpicture}[
  font=\sffamily\footnotesize,
  caja/.style={draw=azulfw,fill=azulclaro,rounded corners=2pt,align=center,minimum height=1.05cm,text width=2.35cm},
  gate/.style={draw=ocrefw,fill=ocreclaro,rounded corners=2pt,align=center,minimum height=1.05cm,text width=2.6cm},
  fin/.style={draw=verdefw,fill=verdeclaro,rounded corners=2pt,align=center,minimum height=1.05cm,text width=2.35cm},
  flecha/.style={-{Stealth[length=2.2mm]},thick,color=grisfw},
  node distance=0.55cm and 0.55cm]
\node[caja] (svc) {Servicio\\\texttt{usuarios}\\(NestJS + Prisma)};
\node[caja,right=of svc] (f3) {Fase 3\\compilación\\(\texttt{npm ci}, \texttt{tsc})};
\node[gate,right=of f3] (g1) {Fase 4: Gate 1\\8 comprobaciones\\(\texttt{gate1.sh})};
\node[caja,right=of g1] (sarif) {5 archivos\\SARIF};
\node[fin,above right=0.25cm and 0.55cm of sarif] (sec) {GitHub\\Security};
\node[caja,below right=0.25cm and 0.55cm of sarif] (motor) {Motor:\\\texttt{framework run}\\\texttt{-{}-fase 4}};
\node[fin,below=0.55cm of motor,text width=3.2cm] (neo) {Neo4j\\Ejecucion → Hallazgo\\→ CWE};
\node[caja,below=0.55cm of g1] (f910) {Fases 9 y 10\\Jest y\\\texttt{docker build}};
\draw[flecha] (svc) -- (f3);
\draw[flecha] (f3) -- (g1);
\draw[flecha] (g1) -- (sarif);
\draw[flecha] (sarif) |- (sec);
\draw[flecha] (sarif) |- (motor);
\draw[flecha] (motor) -- (neo);
\draw[flecha] (g1) -- (f910);
\end{tikzpicture}
\caption{Recorrido del walking skeleton en cada ejecución de GitHub Actions.}
\end{figure}

\section{Fases del pipeline}

\begin{table}[H]
\centering
\caption{Estado de cada fase del workflow al cierre del sprint 1.}
\begin{tabularx}{\textwidth}{@{}P{0.9cm}P{4.2cm}P{2.6cm}L@{}}
\toprule
\textbf{Fase} & \textbf{Nombre} & \textbf{Estado} & \textbf{Detalle}\\
\midrule
1--2 & Requisitos y arquitectura & Pendiente & Sprint 3.\\
3 & Generación y compilación & Versión inicial & Compila el servicio de referencia (20 s).\\
4 & Security Gate 1 & Versión inicial & 8 comprobaciones, SARIF en Security y registro en Neo4j (46 s).\\
5--6 & Refinamiento iterativo & Pendiente & Sprint 5.\\
7 & Integración & Pendiente & Keycloak en el sprint 2; Traefik y RabbitMQ en el sprint 5.\\
8 & Security Gate 2 & Pendiente & Spectral y Conftest el 21 y 22/10.\\
9 & Pruebas automatizadas & Versión inicial & Jest con cobertura (3 s).\\
10 & Contenerización & Versión inicial & \texttt{docker build} de la imagen (31 s).\\
11 & DAST y Security Gate 3 & Pendiente & Trivy image y Syft el 26/10; ZAP y Nuclei el 27/10.\\
\bottomrule
\end{tabularx}
\end{table}

\section{Correcciones del pipeline del 16 de octubre}

En la revisión de la última ejecución aparecieron dos avisos de GitHub y un problema de la caché.
Se corrigen en el pull request de este informe.

\begin{table}[H]
\centering
\caption{Correcciones aplicadas al workflow al cierre del sprint.}
\begin{tabularx}{\textwidth}{@{}P{5.7cm}L@{}}
\toprule
\textbf{Problema} & \textbf{Corrección}\\
\midrule
\texttt{actions/cache@v4} y \texttt{actions/upload-artifact@v5} usan Node.js 20, que GitHub ya no admite. & Se actualizan a \texttt{cache@v5} y \texttt{upload-artifact@v6}, que usan Node.js 24. El único cambio entre versiones es el entorno de ejecución.\\
\texttt{ubuntu-latest} pasa a Ubuntu 26 desde el 19/10/2026. & Se fija \texttt{ubuntu-24.04}, para que el entorno no cambie a mitad del proyecto y los experimentos DSR sean reproducibles.\\
La base de Trivy (unos 130 MB) se guardaba como una entrada de caché nueva en cada ejecución. & La clave pasa a ser diaria: una sola entrada por día, sin agotar el límite de caché del repositorio.\\
\bottomrule
\end{tabularx}
\end{table}

% ======================================================================
\chapter{Incidencias y lecciones aprendidas}
% ======================================================================

Ninguna de las incidencias se resolvió relajando el Gate. Cada una se corrigió en su origen o,
en el caso de los falsos positivos, se documentó con su justificación.

{\small
\begin{longtable}{@{}P{3.3cm}P{4.5cm}P{5.0cm}P{1.2cm}@{}}
\caption{Incidencias del sprint 1 y su corrección.}\\
\toprule
\textbf{Incidencia} & \textbf{Causa} & \textbf{Corrección} & \textbf{PR}\\
\midrule
\endfirsthead
\toprule
\textbf{Incidencia} & \textbf{Causa} & \textbf{Corrección} & \textbf{PR}\\
\midrule
\endhead
Gitleaks marca un secreto en el historial & Token de ejemplo del badge de CircleCI en la plantilla README de NestJS. No es una credencial real. & Se eliminó del README y se registraron sus huellas en \texttt{.gitleaksignore}, sin repetir el valor en ese archivo. & \pr{146}\\
Trivy fs con 2 vulnerabilidades HIGH & Dependencias transitivas del CLI de Prisma: \texttt{deepmerge-ts} (CVE-2026-40345) y \texttt{mysql2} (GHSA-3f6p-5ww8-9rcr). & \texttt{overrides} de npm a las versiones corregidas, tras verificar \texttt{prisma generate}, build, pruebas e imagen. & \pr{146}\\
Hadolint DL3066 & \texttt{USER node} es un nombre; algunos sistemas no pueden comprobar que no sea administrador. & \texttt{USER 1000:1000}, el mismo usuario identificado por número. & \pr{149}\\
Scripts de instalación de npm & Un paquete comprometido puede ejecutar código al instalarse. & \texttt{-{}-ignore-scripts} en CI y Dockerfile, y \texttt{ignore-scripts=true} en \texttt{.npmrc}. & \pr{147}\\
Fallos de herramienta confundidos con hallazgos & Un error de red o de descarga en el runner dejaba el Gate en rojo igual que un hallazgo. & \texttt{gate1.sh} separa ambos casos, reintenta los errores y precarga las imágenes. & \pr{150}\\
SARIF rechazados por GitHub & Las herramientas corren como root en contenedores y pueden dejar archivos vacíos. & El script devuelve la propiedad al usuario del runner y elimina los SARIF vacíos. & \pr{150}\\
Riesgo de \texttt{\$'\textbackslash r'} en bash & Git en Windows convierte los finales de línea a CRLF. & \texttt{.gitattributes} con \texttt{*.sh text eol=lf}. & \pr{150}\\
Workflow en rojo tras fusionar \pr{152} & \texttt{astral-sh/setup-uv@v10} no existe como etiqueta. & Versión fijada \texttt{v10.3.0}. & \pr{153}\\
Trivy se cuelga en local & En Windows recorre \texttt{node\_modules} a través del montaje de Docker y supera los 5 minutos. & Solo ocurre en local. En el runner tarda segundos, así que se valida con el CI o con una copia limpia. & --\\
\bottomrule
\end{longtable}}

\begin{cajaocre}{Lección principal}
Casi todas las incidencias se debieron a diferencias entre el entorno local y el runner: Windows
frente a Linux, finales de línea o versiones no fijadas. Probar cada cambio en un entorno
equivalente al runner (un clon limpio, en Linux, con las mismas imágenes) permite detectarlas
antes de abrir el PR. Fijar versiones de herramientas, acciones y sistema del runner es lo que
hace reproducible el Gate.
\end{cajaocre}

% ======================================================================
\chapter{Seguridad del propio framework}
% ======================================================================

El capítulo 10 del plan exige que el framework se aplique a sí mismo sus principios. Estos son
los controles que ya están activos al cierre del sprint 1:

\begin{table}[H]
\centering
\caption{Controles del capítulo 10 implementados en el sprint 1.}
\begin{tabularx}{\textwidth}{@{}P{4.2cm}L@{}}
\toprule
\textbf{Riesgo} & \textbf{Control activo}\\
\midrule
Scripts de instalación de npm & \texttt{npm ci -{}-ignore-scripts} en CI y en las etapas del Dockerfile; \texttt{.npmrc} del servicio con \texttt{ignore-scripts=true}.\\
Secretos del framework & \texttt{.env} fuera de Git; claves en GitHub Secrets; Gitleaks revisa todo el historial en cada ejecución.\\
Dependencias del framework & Trivy fs revisa \texttt{package-lock.json} (servicio) y \texttt{uv.lock} (motor).\\
Manipulación del conocimiento & Toda la integración pasa por pull request con aprobación del otro integrante y CODEOWNERS; \texttt{semgrep -{}-test} comprueba las reglas propias en cada ejecución.\\
Contenedores & Imagen final sin herramientas de compilación y con usuario no root.\\
Acceso a servicios locales & Los puertos del \texttt{docker-compose.dev.yml} se publican solo en \texttt{127.0.0.1}; el Neo4j de CI vive solo durante la ejecución.\\
Permisos del workflow & \texttt{contents: read} por defecto; \texttt{security-events: write} solo en el trabajo que publica SARIF.\\
\bottomrule
\end{tabularx}
\end{table}

% ======================================================================
\chapter{Planificación del sprint 2}
% ======================================================================

\section{Objetivo e hito}

\begin{cajaazul}{Hito de control del 30 de octubre de 2026}
\emph{Walking skeleton completo (tres gates) con avisos en Telegram.} Es el primer hito importante
del proyecto: no se empieza el sprint 3 sin haberlo cumplido.
\end{cajaazul}

El sprint 2 va del 19 al 30 de octubre: 10 días y 80 horas de equipo. Construye la base de
conocimiento v1 con GraphRAG, las políticas OPA de los tres gates, las versiones iniciales de los
Gates 2 y 3 sobre el servicio de referencia y los avisos de Telegram.

{\small
\begin{longtable}{@{}P{1.55cm}P{5.6cm}P{5.6cm}@{}}
\caption{Cronograma del sprint 2 por integrante.}\\
\toprule
\textbf{Fecha} & \textbf{Marcelo (motor e IA)} & \textbf{Abner (plataforma y seguridad)}\\
\midrule
\endfirsthead
\toprule
\textbf{Fecha} & \textbf{Marcelo (motor e IA)} & \textbf{Abner (plataforma y seguridad)}\\
\midrule
\endhead
19/10 lun & Base de conocimiento: patrones de microservicios en YAML. & Notificador de Telegram: bot y mensaje de progreso editable.\\
20/10 mar & Controles OWASP ASVS y catálogo CWE / OWASP API Top 10. & Notificaciones en cada paso del workflow y aviso de fallo.\\
21/10 mié & Cargador YAML → Neo4j (nodos y relaciones). & Gate 2 inicial: Spectral con reglas OWASP API.\\
22/10 jue & Índice vectorial; fragmentos y embeddings de OWASP/CWE. & Conftest con políticas Rego para Docker Compose.\\
23/10 vie & Consulta GraphRAG (vectores + relaciones). & Realm de Keycloak y guard JWT en el servicio de referencia.\\
26/10 lun & Políticas Rego de los Gates 1 y 2 con \texttt{opa test}. & Gate 3 inicial: Trivy image y Syft (SBOM).\\
27/10 mar & Política Rego del Gate 3 con justificación. & ZAP API scan con límite de tiempo y Nuclei.\\
28/10 mié & Benchmark Neo4j frente a PostgreSQL + pgvector. & Conversor de resultados de ZAP a SARIF.\\
\rowcolor{azulclaro} 29/10 jue & \multicolumn{2}{P{11.6cm}}{Ambos: walking skeleton completo (Gates 1--3 con OPA en CI).}\\
\rowcolor{azulclaro} 30/10 vie & \multicolumn{2}{P{11.6cm}}{Ambos: revisión, demostración con avisos en Telegram, informe del sprint 2 y planificación del sprint 3.}\\
\bottomrule
\end{longtable}}

\section{Dependencias y preparación}

La revisión del sprint 1 deja varias dependencias entre tareas que conviene resolver al empezar
el sprint:

\begin{itemize}
\item \textbf{Telegram:} los secretos \texttt{TELEGRAM\_BOT\_TOKEN} y \texttt{TELEGRAM\_CHAT\_ID} ya
están en GitHub Secrets, así que el notificador puede empezar el 19/10 sin bloqueo.
\item \textbf{Spectral (21/10)} necesita el documento OpenAPI del servicio. Habrá que exportar
\texttt{openapi.json} en CI desde la configuración de Swagger, sin levantar la base de datos.
\item \textbf{Guard JWT (23/10):} el código nuevo de autenticación quedará cubierto
automáticamente por las reglas propias \texttt{jwt-secreto-en-codigo} y
\texttt{jwt-ignora-expiracion} del Gate 1.
\item \textbf{Trivy image y Syft (26/10)} reutilizan la imagen \texttt{fw-usuarios:ci} que ya
construye la fase 10.
\item \textbf{ZAP y Nuclei (27/10)} necesitan el servicio en marcha dentro del runner. Hay que
añadir PostgreSQL como contenedor de servicio en el workflow y aplicar las migraciones
(\texttt{prisma migrate deploy}) antes de arrancarlo.
\item \textbf{Políticas OPA (26--29/10):} \texttt{gate1.sh} ya produce los SARIF y una tabla de
resultados. La política \texttt{gate1.rego} debe tomar esa decisión, y el script pasará a
limitarse a ejecutar las herramientas.
\item \textbf{Duración del CI:} con DAST, el pipeline crecerá desde los 3,5 minutos actuales. El
límite de tiempo de ZAP y las cachés de imágenes deben mantenerlo por debajo de los 30 minutos de
\texttt{timeout} del trabajo.
\end{itemize}

\appendix
% ======================================================================
\chapter{Evidencias}
% ======================================================================

\begin{table}[H]
\centering
\caption{Dónde comprobar cada resultado del sprint.}
\begin{tabularx}{\textwidth}{@{}P{4.0cm}>{\raggedright\arraybackslash\small}X@{}}
\toprule
\textbf{Evidencia} & \textbf{Ubicación}\\
\midrule
Repositorio & \url{https://github.com/PracticasESPE/Framework_Desarrollo_Seguro}\\
Ejecución del hito en \texttt{dev} & \url{https://github.com/PracticasESPE/Framework_Desarrollo_Seguro/actions/runs/38955751601}\\
Pestaña Security (SARIF) & \url{https://github.com/PracticasESPE/Framework_Desarrollo_Seguro/security/code-scanning}\\
Tablero del proyecto & \url{https://github.com/orgs/PracticasESPE/projects/1}\\
Guía de demostración & \texttt{docs/sprints/sprint1-demostracion.md}\\
Script del Gate 1 & \texttt{scripts/gate1.sh}\\
Reglas Semgrep y casos de prueba & \texttt{semgrep\_rules/}\\
Comparación de LLM & \texttt{evaluation/comparacion\_llm/resultados.md}\\
Esquema de Neo4j & \texttt{graph/01\_restricciones.cypher}, \texttt{graph/02\_indices.cypher}\\
\bottomrule
\end{tabularx}
\end{table}

\section*{Cómo reproducir el Gate 1 en local}
Hace falta Docker en marcha y una terminal Linux (o WSL); en Windows, Trivy se vuelve muy lento
al leer el repositorio a través del montaje de Docker. Desde la raíz del repositorio:

\begin{tcolorbox}[colback=white,colframe=grisfw,boxrule=0.4pt,arc=1mm,left=2mm]
\small\ttfamily
cd services/usuarios \&\& npm ci -{}-ignore-scripts \&\& npx prisma generate \&\& cd ../..\\
bash scripts/gate1.sh\\
docker run -{}-rm -v "\$PWD:/src" semgrep/semgrep:1.172.0 semgrep -{}-test semgrep\_rules/
\end{tcolorbox}

\end{document}
